
\documentclass{article}
\usepackage{colortbl}
\usepackage{makecell}
\usepackage{multirow}
\usepackage{supertabular}

\begin{document}

\newcounter{utterance}

\centering \large Interaction Transcript for game `cladder', experiment `full\_v1.5\_default', episode 771 with Qwen3.5{-}4B.
\vspace{24pt}

{ \footnotesize  \setcounter{utterance}{1}
\setlength{\tabcolsep}{0pt}
\begin{supertabular}{c@{$\;$}|p{.15\linewidth}@{}p{.15\linewidth}p{.15\linewidth}p{.15\linewidth}p{.15\linewidth}p{.15\linewidth}}
    \# & \multicolumn{2}{c}{Player} && \multicolumn{2}{c}{Game Master} \\
    \hline

    \theutterance \stepcounter{utterance}  
    & & & \multicolumn{4}{p{0.6\linewidth}}{
        \cellcolor[rgb]{0.9,0.9,0.9}{
            \makecell[{{p{\linewidth}}}]{
                \texttt{\tiny{[P1$\langle$GM]}}
                \texttt{You are an expert in causal inference. The following question is not a typical commonsense query, but rather a meticulously designed question created by a professor specializing in causal inference, intended to assess the students' mastery of the course content.} \\
\texttt{Imagine a self{-}contained, hypothetical world with only the following conditions, and without any unmentioned factors or causal relationships: Gender has a direct effect on occupation and salary. Skill has a direct effect on occupation and salary. Occupation has a direct effect on salary. Skill is unobserved. For individuals who are not male and blue{-}collar workers, the probability of high salary is 26\%. For individuals who are not male and white{-}collar workers, the probability of high salary is 61\%. For individuals who are male and blue{-}collar workers, the probability of high salary is 55\%. For individuals who are male and white{-}collar workers, the probability of high salary is 68\%. For individuals who are not male and with low skill levels, the probability of white{-}collar job is 40\%. For individuals who are not male and with high skill levels, the probability of white{-}collar job is 15\%. For individuals who are male and with low skill levels, the probability of white{-}collar job is 2\%. For individuals who are male and with high skill levels, the probability of white{-}collar job is 24\%. The overall probability of high skill level is 11\%. If we disregard the mediation effect through occupation, would gender negatively affect salary?} \\
\texttt{Start your answer with 'Yes' or 'No', followed by additional reasoning or evidence to support your explanation.} \\
            }
        }
    }
     \\ \\

    \theutterance \stepcounter{utterance}  
    & \multicolumn{4}{p{0.6\linewidth}}{
        \cellcolor[rgb]{0.561,0.737,0.561}{
            \makecell[{{p{\linewidth}}}]{
                \texttt{\tiny{[P1$\rangle$GM]}}
                \texttt{Yes} \\
\\ 
\texttt{**Reasoning:**} \\
\\ 
\texttt{To determine whether gender negatively affects salary in this hypothetical world, we must analyze the causal structure and the specific conditional probabilities provided, strictly adhering to the instruction to disregard the mediation effect through occupation.} \\
\\ 
\texttt{1.  **Causal Structure Analysis**:} \\
\texttt{    The problem states that **Gender** has a direct effect on **Salary**. It also states that **Skill** has a direct effect on **Salary**. While **Occupation** also has a direct effect on **Salary**, the prompt explicitly instructs us to "disregard the mediation effect through occupation." This implies we should focus on the direct causal paths from Gender to Salary, mediated only by Skill (if Skill is a confounder or direct cause) or simply by the direct effect itself.} \\
\\ 
\texttt{    Crucially, the prompt defines the probability of high salary based on specific combinations of Gender and Occupation. However, since Occupation is a mediator (Gender \$\textbackslash{}to\$ Occupation \$\textbackslash{}to\$ Salary) and we are told to disregard this path, we must look at the *net* effect of Gender on Salary after accounting for the distribution of Skill, which is a direct cause of both Occupation and Salary.} \\
\\ 
\texttt{2.  **Analyzing the Data via Skill Levels**:} \\
\texttt{    The prompt provides probabilities of high salary conditional on Gender and Occupation. To find the total effect of Gender on Salary, we need to aggregate these probabilities weighted by the probability of being in each occupation, which is determined by Gender and Skill.} \\
\\ 
\texttt{    Let \$P(S=H)\$ be the probability of high skill = 0.11.} \\
\texttt{    Let \$P(S=L) = 1 {-} 0.11 = 0.89\$.} \\
\\ 
\texttt{    We are given the probability of White{-}Collar (\$W\$) given Gender (\$G\$) and Skill (\$S\$):} \\
\texttt{    *   \$P(W | \textbackslash{}neg M, L) = 0.40\$ \$\textbackslash{}Rightarrow\$ \$P(B | \textbackslash{}neg M, L) = 0.60\$} \\
\texttt{    *   \$P(W | \textbackslash{}neg M, H) = 0.15\$ \$\textbackslash{}Rightarrow\$ \$P(B | \textbackslash{}neg M, H) = 0.85\$} \\
\texttt{    *   \$P(W | M, L) = 0.02\$ \$\textbackslash{}Rightarrow\$ \$P(B | M, L) = 0.98\$} \\
\texttt{    *   \$P(W | M, H) = 0.24\$ \$\textbackslash{}Rightarrow\$ \$P(B | M, H) = 0.76\$} \\
\\ 
\texttt{    Now, let's calculate the probability of High Salary (\$Y=1\$) for Males (\$M\$) and Non{-}Males (\$\textbackslash{}neg M\$) by summing over the skill levels and occupation types.} \\
\\ 
\texttt{    **For Non{-}Males (\$\textbackslash{}neg M\$):**} \\
\texttt{    *   **Low Skill (\$S} \\
            }
        }
    }
    & & \\ \\

    \theutterance \stepcounter{utterance}  
    & & & \multicolumn{2}{p{0.3\linewidth}}{
        \cellcolor[rgb]{0.9,0.9,0.9}{
            \makecell[{{p{\linewidth}}}]{
                \texttt{\tiny{[GM$|$GM]}}
                \texttt{Yes} \\
            }
        }
    }
    & & \\ \\

    \theutterance \stepcounter{utterance}  
    & & & \multicolumn{2}{p{0.3\linewidth}}{
        \cellcolor[rgb]{0.9,0.9,0.9}{
            \makecell[{{p{\linewidth}}}]{
                \texttt{\tiny{[GM$|$GM]}}
                \texttt{no} \\
            }
        }
    }
    & & \\ \\

    \theutterance \stepcounter{utterance}  
    & & & \multicolumn{2}{p{0.3\linewidth}}{
        \cellcolor[rgb]{0.9,0.9,0.9}{
            \makecell[{{p{\linewidth}}}]{
                \texttt{\tiny{[GM$|$GM]}}
                \texttt{game\_result = LOSE} \\
            }
        }
    }
    & & \\ \\

\end{supertabular}
}

\end{document}
